\par\Needspace{6\baselineskip}\medskip
\SourceLeaf{140}\pdfbookmark[2]{100.}{article-100}
\begin{center}100.\end{center}
Bei der Lösung der Gleichung 13 oder $13^*$ betrachten wir zunächst gesondert den Fall, dass $z$ so klein ist, dass sich $Z$ durch eine rasch konvergierende Reihe nach Potenzen von $z$ ausdrücken lässt. Es ergibt sich
\[
\begin{gathered}
(1+2z)\BareRoot{(z+zz)}=z^{\frac12}+\tfrac52z^{\frac32}+\tfrac78z^{\frac52}-\cdots,\\[.5em]
\log(\BareRoot{(1+z)}+\BareRoot z)=z^{\frac12}-\tfrac16z^{\frac32}+\tfrac3{40}z^{\frac52}-\cdots,
\end{gathered}
\]
\noindent Somit ergibt sich als Zähler von $Z=\tfrac83z^{\frac32}+\tfrac45z^{\frac52}-\cdots$ und als Nenner $=2z^{\frac32}+3z^{\frac52}+\cdots$. Daraus folgt $Z=\tfrac43-\tfrac85z+\cdots$. Um das Bildungsgesetz der Reihe zu finden, differenzieren wir die Gleichung
\[
2(z+zz)^{\frac32}Z=(1+2z)\BareRoot{(z+zz)}-\log(\BareRoot{(1+z)}+\BareRoot z),
\]
\noindent woraus sich nach Vereinfachung ergibt
\[
2(z+zz)^{\frac32}\frac{\mathrm dZ}{\mathrm dz}+3Z(1+2z)\BareRoot{(z+zz)}=4\BareRoot{(z+zz)}
\]
\noindent oder
\[
(2z+2zz)\frac{\mathrm dZ}{\mathrm dz}=4-(3+6z)Z,
\]
\noindent woraus man auf ähnliche Weise wie in Art. 90 erhält
\[
\begin{aligned}
Z={}&\tfrac43-\frac{4\mathpunct{.}6}{3\mathpunct{.}5}z+\frac{4\mathpunct{.}6\mathpunct{.}8}{3\mathpunct{.}5\mathpunct{.}7}zz\\[.5em]
&-\frac{4\mathpunct{.}6\mathpunct{.}\overset{\mkern2mu .}{8}\mathpunct{.}10}{3\mathpunct{.}5\mathpunct{.}7\mathpunct{.}9}z^3+\frac{4\mathpunct{.}6\mathpunct{.}8\mathpunct{.}10\mathpunct{.}12}{3\mathpunct{.}5\mathpunct{.}7\mathpunct{.}9\mathpunct{.}11}z^4-\text{usw.}
\end{aligned}
\]
\noindent Offenbar hängt $Z$ also von $-z$ genau so ab wie oben bei der Ellipse $X$ von $x$. Setzen wir daher
\[
Z=\frac1{\tfrac34+\tfrac9{10}(z+\zeta)}\mathpunct{\raisebox{.45ex}{,}}
\]
\noindent so bestimmt sich auch $\zeta$ aus $-z$ ebenso wie oben $\xi$ aus $x$, und es ergibt sich
\[
\mathrm{14)}\qquad\zeta=\cfrac{\tfrac2{35}zz}{1-\tfrac2{35}z+\cfrac{\tfrac{40}{63}z}{1+\cfrac{\tfrac4{99}z}{1+\cfrac{\tfrac{70}{143}z}{1+\text{usw.}}}}}
\]
\noindent oder
\[
\zeta=\cfrac{\tfrac2{35}zz}{1+\tfrac{18}{35}z+\cfrac{\tfrac4{63}z}{1+\cfrac{\tfrac{40}{99}z}{1+\cfrac{\tfrac{18}{143}z}{1+\text{usw.}}}}}
\]
\noindent\SourceLeaf{141}Die Werte für $\zeta$ wurden auf diese Weise im Bereich von $z=0$ bis $z=0{,}3$ in Schritten von einem Tausendstel berechnet. Die dritte Spalte der Tafel III gibt diese Werte an.
